
\begin{obeactivity}
\textbf{Aktivitas 6.1: Analisis Mode ECB vs CBC}
1. Gunakan alat bantu enkripsi (seperti CyberChef) untuk mengenkripsi sebuah gambar sederhana (misal: gambar logo kecil) menggunakan AES mode ECB dan AES mode CBC.
2. Amati hasil cipherteks yang dihasilkan. Apakah gambar asli masih terlihat pada mode ECB?
3. Diskusikan mengapa mode CBC lebih aman untuk data yang memiliki pola berulang.
\end{obeactivity}

Langkah-langkah pada Aktivitas 6.1 dapat dijalankan langsung melalui program
pada Listing~\ref{lst:mode-operasi}. Program itu memakai block cipher mainan
berukuran 4 byte untuk membandingkan ECB dan CBC pada plainteks berpola
berulang, sehingga kebocoran pola pada ECB terlihat jelas tanpa perlu mengolah
gambar.

\lstinputlisting[
    language=Python,
    caption={Mode operasi ECB dan CBC: kebocoran pola pada blok yang berulang},
    label={lst:mode-operasi}
]{\codefile{python/mode_operasi.py}}

Program yang sama tersedia dalam C++ (\texttt{code/cpp/mode\_operasi.cpp}) dan
MATLAB (\texttt{code/matlab/mode\_operasi.m}).
